\section[Python network inference]{Python network inference: preserving the assignment law through learning and policy evaluation}
\label{sec:exec25-python-network}

A bank allocating interventions across connected borrowers needs more than a graph attached to an estimator. A household relationship can make an apparently untreated customer indirectly exposed; a shared employer can create common shocks; a branch allocation can induce correlated assignments; a recommendation policy can select a different population from the one used to train a model. These distinctions determine both the target and the uncertainty. The Python comparison therefore extends the existing MossSpider review with six source collections addressing different parts of the problem: ANI, GraphDML, targeted interference, CausalMotifs, CausalGraphCut and Causal-PIE. None is treated as a universal network-causal engine.~\citep{exec24_interference_mossspider,exec25_pnet_ani,exec25_pnet_dml,exec25_pnet_targeted,exec25_pnet_motifs,exec25_pnet_cut,exec25_pnet_pie}

The selected originals are pinned to commits dated 8 June 2025, 7 October 2026, 10 August 2026, 8 February 2021, 26 May 2025 and 21 March 2025 respectively. GraphDML declares version 0.1.0. Its selected commit is at 15:30 UTC on the cutoff day, or 23:30 in Asia/Shanghai, within the client cutoff. Forty-eight saved files, complete selected repository trees and inspected technical anchors support the comparison. Repository search supplies a bounded candidate list with 54 unique leads; uninspected leads and search truncation remain omissions. The calculations below are project-authored checks of specified expressions. No acquired package, upstream test, neural training routine or graph simulator was executed, and papers mentioned in source comments are not silently adopted as theorem evidence.~\citep{exec25_pnet_choice}

\subsection{A common interface must preserve distinct network targets}

Write $D_i$ for the treatment assigned to node $i$, $A$ for the declared network, and $e_i(D,A)$ for its exposure. The response may be represented as $Y_i(d,e)$ only under a substantive restriction connecting the full assignment vector to the chosen exposure. A graph feature learned for prediction does not establish that restriction. A design-based contrast requires assignment probabilities for exposure events; a partially linear model instead imposes a response structure; a policy evaluator needs the law of the proposed assignments and the relevant time horizon. These objects should not be collapsed into a generic effect field.

For banking use, the request should therefore bind the observation time, node identifiers, arrow direction, edge weights, treatment assignment law, exposure operator, analysis population, target policy and uncertainty law. A sum of neighbors' treatments measures a different intervention from their mean, especially when degree varies. An isolated borrower needs an explicit exposure convention. A missing relationship is not automatically an absent relationship. The proposed CLI and MCP result must carry these decisions and any refusal, so downstream agents cannot interpret an output number without its scientific contract. The package comparisons provide concrete tests of this proposed interface; they do not establish that the interface is implemented.~\citep{exec25_pnet_choice}

\subsection{ANI: uncertainty must not depend on the order of requested contrasts}

ANI constructs an inverse-probability exposure contrast from supplied indicators and probabilities. Its population mask selects the contributors, while excluded score entries receive a sentinel that is not intended for estimation. The standard-error routine computes distances in the full undirected, unweighted graph and assumes nodes are numbered consistently with score positions. Thus subsetting outcomes does not silently mean recomputing distances on an induced graph. The explicit zero-bandwidth branch uses the population-standard-deviation convention divided by the square root of the selected count.~\citep{exec25_pnet_ani}

For a positive bandwidth, let $z$ be the centered score vector over the selected population of size $n$ and let $K_{ij}=1\{d_A(i,j)\leq b\}$. The selected variance of the mean is $z'Kz/n^2$. When the intermediate quadratic is nonpositive, the code constructs a replacement from overlapping neighborhoods. If $B_{ij}=1\{d_A(i,j)\leq b\}$ and $r_i=\sum_jB_{ij}$, that replacement is $P=\widetilde B\widetilde B'$ with $\widetilde B_{ij}=B_{ij}/\sqrt{r_i}$. Its positive-semidefinite form is apparent from the product, but the implementation places the replacement into a weights variable shared across the loop over contrasts. Subsequent contrasts use that replacement even when their original quadratic would have been positive. The fallback flag records whether the current contrast triggered the branch, not necessarily which kernel it inherited.~\citep{exec25_pnet_ani}

An authored three-node path illustrates the interface risk. Set $b=1$ and
\[
 K=\begin{pmatrix}1&1&0\\1&1&1\\0&1&1\end{pmatrix},
 \qquad u=(1,-2,1)',\qquad v=(1,0,-1)'.
\]
Both scores are centered. Their original variance quadratics divided by nine are $-2/9$ and $2/9$. Processing $u$ first activates the overlap kernel. The resulting variance for $u$ is $(7-16/\sqrt6)/9$, and $v$ then receives $1/9$. Processing $v$ first instead gives it $2/9$. The latter contrast can therefore change solely because another requested contrast preceded it. This is a source-expression witness, not a reported run of ANI on a fitted banking sample.~\citep{exec25_pnet_choice}

ANI is a conditional reuse candidate for design-based network uncertainty. The adapter should preserve the original kernel, allocate any fallback locally to each contrast, report the kernel actually used, and test permutation invariance of a batch of contrasts. It must also test disconnected and edgeless graphs, singleton components, node ordering, structural support, nonfinite scores and subsampling. These engineering tests do not prove the sampling assumptions needed for a network central-limit approximation. The authentic BSD three-clause license is a useful starting point; installed dependencies and distribution obligations remain a separate review.~\citep{exec25_pnet_ani,exec25_pnet_choice}

\subsection{GraphDML: useful graph-aware construction with population and covariance gates}

GraphDML makes several interface distinctions explicit. The data container checks covariate finiteness and graph dimensions, warns when dropping self-loops, rejects a nonsymmetric graph unless direction is declared, and checks node identifiers. Its exposure operators distinguish unweighted sums, unweighted means, weighted sums and means, and shortest-path distance exactly two. The two-hop implementation removes self and direct-neighbor entries from the thresholded matrix product. The mean operator leaves isolated rows at zero. Those conventions are valuable adapter contracts, provided the bank accepts their causal meaning.~\citep{exec25_pnet_dml}

The estimation path forms dependency sets from own and exposure-row support. For a held-out focal set, buffered nuisance training excludes treatment observations in the union of its dependency sets and excludes outcome observations whose own dependency sets overlap that union. The default final covariance is an empirical sandwich without a supplied group vector. This is not evidence that the package lacks a cluster option: a separate cluster-summary method aggregates scores by a supplied group and explicitly requires one repetition. Disjoint graph supports alone do not verify conditional independence of primitive disturbances. A common employer shock can violate that interpretation despite distant graph nodes; the scientific model and a justified cluster law remain necessary.~\citep{exec25_pnet_dml}

The population contract also needs attention. When the focal set is a named selection rule, the declared estimation nodes are supplied as candidates. When explicit focal indices are supplied, the selected branch checks index range and disjoint dependency supports but returns the indices without intersecting them with the declared estimation population. Total-effect weights subsequently use that declared population. An adapter should reject incompatible index sets or explicitly justify transport of the fitted coefficients; it should not silently describe coefficients from one population as estimated in another. This is a branch-level finding, not a claim that every fit takes the incompatible path.~\citep{exec25_pnet_dml}

The total effect has a precise model interpretation. For the proposed partially linear response
\[
 Y_i=g(X_i)+\theta D_i+\sum_{k=1}^q\beta_k(E_kD)_i+\varepsilon_i,
\]
with appropriately identified common coefficients, changing every treatment from zero to one gives $\theta+\sum_k\beta_k\overline{E_k\mathbf1}$ over the declared population. This follows by subtraction under the stated model; it does not identify an unrestricted heterogeneous policy effect. The source computes these weights and uses their covariance quadratic for total-effect uncertainty.~\citep{exec25_pnet_dml,exec25_pnet_choice}

For repeated fits, the selected helper takes a coordinatewise median of coefficients and then an elementwise median of covariance matrices adjusted by coefficient-displacement outer products. Elementwise medians need not preserve positive semidefiniteness. Consider identical coefficient vectors and three positive-definite covariance inputs
\[
 C_1=\begin{pmatrix}1&-2\\-2&5\end{pmatrix},\quad
 C_2=\begin{pmatrix}5&-2\\-2&1\end{pmatrix},\quad
 C_3=\begin{pmatrix}1&0\\0&1\end{pmatrix}.
\]
Each has positive leading entry and determinant one. Their elementwise median is $M=\left(\begin{smallmatrix}1&-2\\-2&1\end{smallmatrix}\right)$, with determinant $-3$. For total-effect weights $(1,1)$, the variance quadratic is $-2$. This exact helper-input counterexample does not establish that a particular fitted repetition produces these matrices. It establishes that entrywise aggregation alone cannot certify a covariance suitable for every linear contrast. Merely clipping a negative variance does not establish correct coverage.~\citep{exec25_pnet_choice}

The proposed decision is conditional adoption through a narrow adapter: retain graph-aware splits, exposure operators, diagnostics and explicit cluster summaries, while requiring population compatibility and contrast-valid covariance aggregation. Public trials should include common shocks, near-singular moments, repeated splitting, isolated nodes, directed graphs, exact two-hop semantics and changed node ordering. The MIT original and dependency declarations permit a concrete rights review; they are not complete artifact clearance. No native capability is admitted from these static findings.~\citep{exec25_pnet_dml,exec25_pnet_choice}

\subsection{Targeted interference and CausalMotifs: learning the representation is not the entire inference problem}

The targeted-interference implementation combines a graph encoder, separate outcome heads, a treatment-probability head, a neighborhood-density head and spline fluctuations. The selected neighborhood mean is $AD/(A\mathbf1)$ with no isolate guard in that expression. The treatment head is trained with binary cross-entropy, while the fluctuation correction uses the product of its output and the neighborhood-density output plus a small numerical stabilizer. The treatment of control targets and the meaning of that factorization need a method-specific derivation and known-law comparison. The inspected expression alone does not justify declaring the entire method incorrect.~\citep{exec25_pnet_targeted}

The selected code also includes CUDA-specific allocations and evaluates supplied test potential-outcome truths during training methods. The inspected top-level loop uses configured epoch counts; it does not establish that the optimizer selects a checkpoint using test truth. Nevertheless, a proposed public experiment should keep final evaluation oracles inaccessible until the model and hyperparameters are frozen. Synthetic potential outcomes can supply that oracle, whereas observed bank data generally cannot. The conditional recommendation is to study the architecture as a research reference and defer direct adoption until the target derivation, isolate handling, device portability, held-out evaluation and distribution rights are resolved. No license-named original appears in the selected complete tree.~\citep{exec25_pnet_targeted,exec25_pnet_choice}

CausalMotifs learns partitions using network-motif features and simulated exposure probabilities. Its constructor randomly divides rows into training and estimation subsets. That split is useful for separating partition selection from estimation on rows, but does not itself make their network disturbances independent. The nonseparate validation branch removes rows with zero empirical probabilities and fits weighted least squares using the default fit covariance. The separate H\'ajek helper contains individual variance and covariance terms rather than explicit pairwise exposure covariances. Appropriate uncertainty therefore needs a declared assignment/dependence law and a target population that records the filtering.~\citep{exec25_pnet_motifs}

Two selected support checks make small, deterministic tests especially valuable. The branch without an own-treatment restriction reads the stored probability collection rather than the passed collection. The own-treatment branch adds a list of rule arrays to the treatment array before comparing with the target value; it does not append a parenthesized treatment-condition array to the rule list. In the scalar arithmetic analogue, a true rule and treatment one produce $(1+1)==1$, which is false although their conjunction is true. Separately, the caller computes left-control and right-control support checks but repeats the right-control term in its rejection predicate. Array shapes and the exact caller path must be qualified in native tests before interpreting the practical effect.~\citep{exec25_pnet_motifs,exec25_pnet_choice}

The example generator requires equal care. It uses a very large synthetic graph, clustered labels, and an index-distance shortcut to judge selected neighbor-neighbor relationships. Its subsequent probability simulation rotates cluster alignment, unlike the initial assignment expression. A bank should not port those illustrative shortcuts as general graph or assignment semantics. The MIT source is a conditional research/adapter candidate after explicit Boolean-rule tests, assignment-law reconciliation, support accounting and dependent uncertainty. Native qualification remains open.~\citep{exec25_pnet_motifs,exec25_pnet_choice}

\subsection{CausalGraphCut: inspect the custom delegate before judging the interface}

CausalGraphCut addresses experimental allocation rather than simply estimating an effect after assignments have been made. It constructs a covariance-weighted graph objective and delegates spectral partitioning to a local class. That class directly eigendecomposes its input for the selected discretization path. It would therefore be wrong to infer a standard-library affinity-interface error merely from the public class name: the actual delegate has different semantics. The proposed reuse is an allocation-design research reference whose objective must match the bank's assignment probability, network and outcome dependence.~\citep{exec25_pnet_cut}

The inspected source exposes separate implementation gates. The default maximum cluster count is a rounded power of the node count cast to a signed eight-bit integer. For 4,096 nodes the unwrapped value is 256, whose eight-bit wrapping representation is zero. This is an authored arithmetic witness; installed NumPy conversion behavior remains a native test. The online routine resets batch size to five, uses complete batches only, and depends on a graph variable populated by the main example rather than an explicit graph argument. With seven requested samples, its batch arithmetic schedules five and omits two. With only one complete batch, the returned cluster variable is not assigned in the initial branch. Larger studies also need to retain the full allocation and adaptation history.~\citep{exec25_pnet_cut,exec25_pnet_choice}

The objective list is sorted in descending order and the online caller selects its last element, which correctly corresponds to the smallest listed objective; this is not a sorting defect. Other scope distinctions remain: the selected objective helpers use half-probability constants, while design classes accept a probability parameter and use different neighborhood or cluster probability exponents. Generalizing the objective to arbitrary assignment probabilities requires a derivation, not just changing a parameter. Direct adoption is deferred pending explicit graph arguments, safe integer handling, sample accounting, typed return paths, adaptive-design validation and rights. No license-named original appears in the selected complete tree.~\citep{exec25_pnet_cut,exec25_pnet_choice}

\subsection{Causal-PIE: an embedding match is not the assignment probability}

Causal-PIE connects learned neighborhood representations to policy evaluation. In its one-stage path, records are selected when own action matches the target and a learned embedding is sufficiently close. The relative criterion uses a strict inequality between root-mean-square distance and a threshold times the root-mean-square reference embedding. Identical zero embeddings fail this test because $0<0$ is false. Moreover, filtering and probability calculation reverse which embedding supplies the reference scale. An adapter must specify a symmetric or intentionally asymmetric equivalence rule, zero-scale behavior and the scientific adequacy of the learned exposure representation.~\citep{exec25_pnet_pie}

The probability calculation offers Monte Carlo and traversal modes. Monte Carlo draws independent neighbor actions using the supplied Bernoulli probability $p$. Traversal enumerates binary configurations but its active expression weights them uniformly. After clipping the neighborhood result, the code multiplies by $p$ regardless of whether the target own action is zero or one. These choices coincide with some balanced designs, but not arbitrary assignment laws.~\citep{exec25_pnet_pie}

For a transparent authored example, suppose one neighbor is treated independently with $p=1/5$, the target own action is zero, and the target neighbor action is one. Assume a nonzero injective embedding, no added noise and an inactive clipping bound. The correct joint assignment probability is $(1-p)p=4/25$. Uniform traversal followed by multiplication by $p$ gives $1/10$; the Monte Carlo neighborhood limit followed by that same own-action factor gives $1/25$. The experiment should test both probability modes at balanced and unbalanced assignments and for both own actions. A floor on the estimated probability cannot repair the wrong assignment law.~\citep{exec25_pnet_choice}

A second composed-path witness concerns empty support. When no record passes the filter, the selected routine returns record index zero with probability one. The caller then computes a weighted residual correction using the returned record and divides by the full number of node-time observations. If that first residual were six and the full count ten, the correction would be $3/5$ rather than the zero sum over an empty selection. This is hypothetical arithmetic, not a measured dataset. The appropriate product response should report empty selected support, and either return the explicitly defined plugin-only quantity with its limitations or refuse the requested correction. A numerical sentinel must not become an observed policy match.~\citep{exec25_pnet_pie,exec25_pnet_choice}

The multi-stage implementation poses different questions. It combines a fitted value function, an RKHS-trained density ratio, discounted next values and an initial-value plugin normalized by $1-\gamma$. Its selected constructor assigns the same trajectory indices to training and test collections. The residual uses the node, trajectory and $T-1$ transition count, so finite-horizon terminal handling and stationary-policy interpretation need separate qualification. The RKHS bandwidth is a median embedding distance, creating a zero-bandwidth gate when embeddings collapse. These are concrete proposed tests; neither the source's double-robust naming nor an accurate fitted reward model establishes valid policy uncertainty.~\citep{exec25_pnet_pie}

The candidate is retained as a policy-evaluation research reference, with direct adoption deferred until probability-law weighting, empty support, matching scale, trajectory separation, horizon semantics, collapsed representations and rights are resolved. The selected complete tree contains no license-named original. Public synthetic trajectories with known intervention values can test these contracts before any expensive bank study. Unknown confounding and an inadequate exposure map still require scientific restrictions or more information.~\citep{exec25_pnet_choice}

\subsection{Integration decisions and a bounded public qualification sequence}

The six decisions are independent. ANI and GraphDML supply concrete adapter candidates after uncertainty and population repairs; CausalMotifs supplies a conditional partition-learning reference after support and dependence checks. The neural targeted, allocation-design and policy-evaluation repositories supply useful architecture and experiment references, with direct adoption deferred until their stated rights and technical gates are resolved. A port to another language should preserve the declared mathematical contract, not reproduce an accidental sentinel or change the population silently. Direct reuse is preferred only when the actual artifact, license and scientific regime meet that contract.~\citep{exec25_pnet_choice}

The first public qualification stage should enumerate tiny assignment spaces and test exposure orientation, isolate behavior, exact joint support, own-action probabilities and policy denominators. The second should examine covariance aggregation, batch-order invariance, common shocks and graph-aware splitting on frozen known-law data. The third should freeze model selection and evaluate all requests, including failures, empty support and nonfinite outputs, on held-out graph and trajectory families. A comparison records identification assumptions, sampling and Monte Carlo uncertainty, runtime, complete costs, refusal accuracy and target-population coverage separately. It cannot manufacture missing bank assignment truth or prove universal validity under arbitrary hidden network dependence.~\citep{exec25_pnet_choice}

For scale, a hypothetical USD30,000 study costs USD30 per 1,000 attempted requests, USD50 per 600 completed requests or USD75 per 400 scientifically accepted outputs. These are three distinct denominators, not observed return on investment. Bounded search tails, complete dependency/artifact review, native reproduction and independent acceptance remain open. The selected primary sources and exact authored witnesses establish a reviewable route to public qualification, not complete Python network qualification.~\citep{exec25_pnet_choice}
